\begin{figure}[t]
\centering
\begin{tikzpicture}[
  font=\footnotesize, >={Stealth[length=1.5mm]},
  io/.style={draw, rounded corners=2pt, minimum width=3.9cm, minimum height=9mm, align=center},
  stp/.style={draw, rounded corners=1pt, fill=white, minimum width=2.45cm, minimum height=9mm, align=center},
  dot/.style={circle, fill=#1, inner sep=0pt, minimum size=1.6mm},
  job/.style={fill=#1, minimum size=1.7mm, inner sep=0pt}
]
% ---- inputs
\node[io, fill=red!10]  (wf)   at (2.0,4.4) {\hspace{9mm}Scientific workflow\\\hspace{9mm}+ time-to-result};
\node[io, fill=gray!10] (hist) at (6.6,4.4) {\hspace{7mm}Historical\\\hspace{7mm}executions};
% mini DAG icon
\begin{scope}[shift={(0.45,4.4)}]
  \node[dot=teal]   (a) at (0,0.22) {};
  \node[dot=violet] (b) at (-0.22,0) {}; \node[dot=teal] (c) at (0.22,0) {};
  \node[dot=orange] (d) at (0,-0.22) {};
  \draw (a)--(b) (a)--(c) (b)--(d) (c)--(d);
\end{scope}
% database icon
\begin{scope}[shift={(5.15,4.4)}]
  \draw[fill=white] (-0.2,0.2) -- (-0.2,-0.2) arc[start angle=180, end angle=360, x radius=0.2, y radius=0.07] -- (0.2,0.2);
  \draw[fill=white] (0,0.2) ellipse[x radius=0.2, y radius=0.07];
  \draw (-0.2,0) arc[start angle=180, end angle=360, x radius=0.2, y radius=0.07];
\end{scope}

% ---- SOU
\node[draw, rounded corners=2pt, fill=gray!4, minimum width=8.6cm, minimum height=21mm] (sou) at (4.3,2.45) {};
\node[anchor=north, font=\footnotesize\bfseries, inner sep=2pt] at (sou.north) {SOU};
\node[stp, fill=blue!8]                (pred) at (1.4,2.6)  {Runtime\\prediction};
\node[stp, fill=blue!8, very thick]    (sim)  at (4.3,2.6)  {Makespan\\simulation};
\node[stp, fill=orange!15, dashed]     (crit) at (7.2,2.6)  {Offloading criterion\\(pluggable)};
\draw[->] (pred) -- (sim);
\draw[->] (crit) -- (sim);
\draw[->] (sim.south) ++(0.35,0) arc[start angle=0, end angle=-180, x radius=0.35, y radius=0.2];
\node[font=\scriptsize\itshape, anchor=north, inner sep=1pt] at (4.3,1.93) {until time-to-result is met};

% ---- clusters
\node[io, fill=green!8] (c1) at (2.0,0.55) {\hspace{9mm}On-premises\\\hspace{9mm}cluster};
\node[io, fill=green!8] (c2) at (6.6,0.55) {\hspace{9mm}Shared HPC\\\hspace{9mm}or cloud};
\foreach \x/\y/\c in {0/0/teal,1/0/violet,2/0/teal,0/1/orange,1/1/teal,2/1/violet}
  \node[job=\c] at ($(0.3,0.45)+(\x*0.23,\y*0.23)$) {};
\foreach \x/\y/\c in {0/0/violet,1/0/teal,0/1/teal}
  \node[job=\c] at ($(4.9,0.45)+(\x*0.23,\y*0.23)$) {};

% ---- flow
\draw[->] (wf)   -- (wf |- sou.north);
\draw[->] (hist) -- (hist |- sou.north);
\draw[->] (c1 |- sou.south) -- node[font=\scriptsize\itshape, right] {local jobs} (c1);
\draw[->] (c2 |- sou.south) -- node[font=\scriptsize\itshape, right] {offloaded jobs} (c2);
\end{tikzpicture}
\caption{SOU overview: predicted job runtimes drive a makespan simulation in which a pluggable offloading criterion selects jobs for the secondary cluster until the user-defined time-to-result is met.}
\label{fig:intro}
\end{figure}


\paragraph{Workflow Makespan Prediction}
SOU operates by first predicting job runtimes from historical executions~\cite{bader_lotaru_2024, bader_sizey_2024, balis_evaluation_2023} to simulate the federated makespan prior to submission.
For each task, we collect the runtimes $\mathcal T$ of its historical jobs and their input size $x$. We define $x$ as the total size of the initial workflow inputs that a job and its ancestors consume, which we obtain by reconstructing the job DAG.
If $x$ and the runtime show a Pearson correlation of $\rho\geq0.8$, we fit a task-specific model with intercept $\beta_0$ and coefficient $\beta_1$, either by linear regression $t=\beta_0+\beta_1x$ or by Bayesian ridge regression $t\sim\mathcal N(\beta_0+\beta_1x,\sigma^2)$ to additionally capture uncertainty. Thus,
\[
\hat t_q(j)=
\begin{cases}
\mu_j+\Phi^{-1}(q)\,\sigma_j, & \rho\geq0.8,\\
\operatorname{median}(\mathcal T), & \text{otherwise}.
\end{cases}
\]
For the Bayesian model, $\mathcal N(\mu_j,\sigma_j^2)$ is the posterior predictive distribution, and quantiles $q>0.5$ give more conservative predictions that serve as a lever to meet the time-to-result with higher assurance~\cite{park_3sigma_2018}.

\paragraph{Heuristic Job Offloading Simulation}
% Algorithm
\begin{algorithm}[t]
\caption{Job Offloading via Makespan Prediction}
\label{alg:offloading}
\small
\SetKwFunction{Sim}{Simulate}
\SetKwProg{Fn}{Function}{:}{}
\KwIn{jobs $J$, target makespan $D$, criterion $\mathcal{C}$}
\KwOut{offloaded jobs $J_\mathrm{off}$}
$J_\mathrm{off}\gets\emptyset$\;
$(\hat M, J_\mathrm{def})\gets\Sim{$J, J_\mathrm{off}$}$\;
\While{$\hat M > D$ \textbf{and} $J_\mathrm{def}\setminus J_\mathrm{off}\neq\emptyset$}{
  add the first job of $J_\mathrm{def}\setminus J_\mathrm{off}$ selected by $\mathcal{C}$ to $J_\mathrm{off}$\;
  $(\hat M, J_\mathrm{def})\gets\Sim{$J, J_\mathrm{off}$}$\;
}
\Return the $J_\mathrm{off}$ with the lowest $\hat M$ found\;
\BlankLine
\Fn{\Sim{$J, J_\mathrm{off}$}}{
  $\hat M\gets 0$;\ $J_\mathrm{def}\gets\emptyset$\;
  \While{unfinished jobs in $J$}{
    \ForEach{ready job $j$}{
      $c\gets$ secondary \textbf{if} $j\in J_\mathrm{off}$ \textbf{else} primary\;
      \lIf{resources $\mathbf r(j)$ are free on $c$}{run $j$ on $c$ for $\hat t_(j)$}
      \lElse{defer $j$;\ $J_\mathrm{def}\gets J_\mathrm{def}\cup\{j\}$}
    }
    advance $\hat M$ to the next job completion\;
  }
  \Return $(\hat M, J_\mathrm{def})$\;
}
\end{algorithm}
SOU selects jobs for offloading by simulating the workflow execution~\cite{casanova_wrench_2020} with the predicted runtimes (Algorithm~\ref{alg:offloading}). It takes the jobs $J$, the time-to-result target $D$, and a selection criterion $\mathcal{C}$ and returns the set of jobs to offload, as a greedy alternative to searching all $2^{|J|}$ candidate sets.

Starting without offloading, it simulates the workflow once (lines 1--2). As long as the predicted makespan $\hat M$ exceeds $D$, it selects one more job by $\mathcal{C}$, adds it to $J_\mathrm{off}$, and re-simulates with this job on the secondary cluster (lines 3--5).
\emph{Longest Job First} (LJF) selects the job with the longest predicted runtime, as long-running jobs best amortize offloading overheads such as data transfer.
As illustrated in Figure~\ref{fig:approach}, LJF first offloads $j_2$ (30\,min), which lowers the predicted makespan from 73 to 43\,min. As this exceeds the 40\,min target, it offloads the next-longest job ($j_4$, 20\,min) and reaches 38\,min. \emph{Smallest Input Size First} (SISF) selects the job with the smallest input size $x$ and targets settings where offloading cost grows with data volume. Further criteria can be implemented.
In contrast, \emph{Primary Environment Fully Occupied} (PEFO) naively offloads whenever the primary cluster is saturated.

To predict the makespan of each selection, \code{Simulate} advances jobs by their predicted runtimes, running jobs in $J_\mathrm{off}$ on the secondary and all other jobs on the primary cluster (lines 11--12). A job whose resources $\mathbf r(j)$ are not free on its cluster waits and is recorded in $J_\mathrm{def}$ (line 13). The function returns $\hat M$ and $J_\mathrm{def}$ (line 15). Candidates for the next selection are those that had to wait for resources in the previous simulation, $J_\mathrm{def}\setminus J_\mathrm{off}$. Once $D$ is met or no candidate remains, SOU returns the configuration with the lowest predicted makespan (line 6).

\begin{figure}[t]
\centering
\begin{tikzpicture}[
  font=\footnotesize, >={Stealth[length=1.3mm]},
  job/.style={circle, draw=gray!70!black, fill=white, minimum size=5.2mm, inner sep=0pt, font=\tiny},
  off/.style={job, fill=orange!35, draw=orange!80!black},
  new/.style={off, very thick},
  edge/.style={->, gray!70, >={Stealth[length=0.9mm]}},
  nxt/.style={->, thick, gray!60!black}
]
% one snapshot: #1 x-centre, #2-#4 styles of the three middle jobs
\newcommand{\snapdag}[4]{%
  \begin{scope}[shift={(#1,1.45)}]
    \node[job] (s) at (0,0.75) {5};
    \node[#2]  (a) at (-0.62,0) {30};
    \node[#3]  (b) at (0,0) {12};
    \node[#4]  (c) at (0.62,0) {20};
    \node[job] (t) at (0,-0.75) {8};
    \foreach \m in {a,b,c} {\draw[edge] (s) -- (\m); \draw[edge] (\m) -- (t);}
  \end{scope}}
% makespan bar: #1 x-centre, #2 length, #3 colour, #4 label
\newcommand{\mkbar}[4]{%
  \begin{scope}[shift={({#1-1.1},0.2)}]
    \fill[#3] (0,-0.07) rectangle (#2,0.07);
    \draw[red!70!black, dashed, thick] (1.6,-0.17) -- (1.6,0.2);
    \node[font=\tiny, anchor=north west, inner sep=1pt] at (0,-0.08) {#4};
  \end{scope}}

% strategy strip
\node[draw, rounded corners=2pt, fill=gray!5, minimum width=8.6cm, minimum height=5.5mm] (strat) at (4.3,2.85) {};
\node[anchor=west, font=\scriptsize] at (strat.west) {\textbf{Offloading Criterion:}};
\node[draw=orange!80!black, fill=orange!20, rounded corners=1pt, font=\scriptsize, inner sep=1.5pt] (ljf) at (3.25,2.85) {LJF};
\node[font=\scriptsize, text=gray, anchor=west] at ($(ljf.east)+(0.05,0)$) {SISF $\mid$ PEFO};
\node[font=\scriptsize, anchor=east] at ($(strat.east)+(-0.1,0)$) {offloading order: 30 $\succ$ 20 $\succ$ 12};

\snapdag{1.25}{job}{job}{job}
\snapdag{4.3}{new}{job}{job}
\snapdag{7.35}{off}{job}{new}

\mkbar{1.25}{2.2}{red!45}{makespan 70 $>$ 50}
\mkbar{4.3}{1.75}{red!45}{makespan 55 $>$ 50}
\mkbar{7.35}{1.45}{teal!55}{makespan 45 $\le$ 50}

\foreach \x/\t in {1.25/start, 4.3/iteration 1, 7.35/iteration 2} \node[font=\tiny\itshape, text=gray!80!black] at (\x-0.95,2.3) {\t};
\draw[nxt] (2.2,1.45) -- node[font=\tiny, above] {offload} (3.35,1.45);
\draw[nxt] (5.25,1.45) -- node[font=\tiny, above] {offload} (6.4,1.45);

\node[font=\tiny, red!70!black, anchor=south, inner sep=1pt] at (1.75,0.4) {objective};
\node[font=\footnotesize, teal!70!black, anchor=west] at (7.92,0.2) {\ding{51}};
\end{tikzpicture}
\caption{Job selection: starting from a DAG annotated with predicted runtimes, the pluggable strategy (here LJF) offloads the next-ranked job and re-simulates the makespan until the objective (here time-to-result) is met.}
\label{fig:approach}
\end{figure}





